Prefix caching reuses prefill only across an exactly shared prefix, so one changed
field invalidates the entire downstream cache.
Yet overwriting the field's own key/value vectors and reusing the rest leaves the model acting on
the \emph{old} value. The reason, established causally across four model families: at prefill the model
has already written the \emph{field-conditioned conclusion} onto downstream
\emph{notes}; the field's own key/value drives under $1\%$ of the decision. Read as a
notebook of memoized conclusions, two capabilities follow.
\textbf{It is editable.} A salient \emph{erratum} amends the notes;
and with \emph{chain-of-thought}---not scale---editing the field alone recovers the decision ($1.00$ at
$8$B, ${\sim}1\%$ compute), while without CoT it is ignored.
\textbf{It is composable.} The notes are position-portable, so a precompiled skill can be RoPE-repositioned
and spliced into any context---indistinguishable from full recompute (logit cosine $0.90$--$0.999$, twelve
models) at $O(L)$ rather than $O(L^2)$ time-to-first-token. A unified edit+compose agent stays
decision-identical to recompute at up to $14.9\times$ lower latency.
The approach applies to any per-token attention KV cache, validated across scale, quantization,
Mixture-of-Experts, and multimodal caches, and extends to several attention variants through small
adapters. Because the erratum is append-only, it composes with production prefix caching: in an online
vLLM benchmark it keeps the prefix cache-aligned ($98.5\%$ hit-rate), cutting $p90$ time-to-first-token by
$53$--$398\times$.
